\section{Our Two-Server PIR Scheme}\label{sec:two-server}

\begin{figure}[p]
    \begin{minipage}{\textwidth}
        \begin{mdframed}
            \label{alg:one-query-ppwps-two}
            \begin{center}
                {\bf Two-server scheme for $Q = \sqrt{n} \log \kappa \cdot \alpha$ queries}
            \end{center}

            \parhead{Offline preprocessing.} 
            %\ \\[-6pt]
%{\it Client}:
            \begin{itemize}[leftmargin=3mm,itemsep=1pt,topsep=1pt]
            %    \item 
                            %, where $M_1=\sqrt{n}\log \kappa \cdot \alpha(\kappa)$. , such
                %that $\alpha(n)=\omega(1)$. 
                % Initialize the parities $p_1,\dots,p_{M_1}$ to zeros.
                
                
                \item 
{\it Hint table.}
Let $M_1=\sqrt{n}\log \kappa \cdot \alpha(\kappa)$. For each $i\in [M_1]$, 
sample a fresh \ppps key $\sk_i$, 
send $\sk_i$ 
to the \underline{right server} 
and receive a parity $p_i := \oplus_{j \in \Set(\sk_i)} \DB[j]$ back. 
%and receive their parities $p_1,\ldots,p_{M_1}$. 
Let $T := \{(\psk_i, p_i )\}_{i \in
                [M_1]}$ denote the client's {\it hint table}.
                \item 
{\it Replacement entries.}
For each chunk $\ell \in \{0,\ldots,\sqrt{n}-1\}$,
repeat the following 
$M_2=3\log \kappa \cdot \alpha(\kappa)$  times:
sample a random index $r_1 \in \{0, \ldots, n-1\}$ in chunk $\ell$, 
send $r_1$ to the \underline{left server}, and receive $\DB[r_1]$.
Store the tuple $(r_1,\DB[r_1])$. 

Similarly, 
for each chunk $\ell \in \{0,\ldots,\sqrt{n}-1\}$,
repeat the following $M_2$  times:
sample a random index $r_2$ in chunk $\ell$, 
send $r_2$ to the \underline{right server}, and receive $\DB[r_2]$.
Store the tuple $(r_2,\DB[r_2])$. 
            \end{itemize}

%\vspace{5pt}
%{\it Right server}:
%upon receiving a \ppps key $\sk$, return $\oplus_{i \in \Set(\sk)}\DB[i]$.

            
            \parhead{Query for index $x \in \{0, 1, \ldots, n-1\}$.}
            
            \begin{enumerate}[leftmargin=5mm,itemsep=1pt]
                \item 
                {\bf Step 1: (Client Querying)}

                \begin{itemize}[leftmargin=5mm,itemsep=1pt]
                        \item Find the first entry $(\sk,p)$ in the hint table such that $x\in \Set(\sk)$.~\footnote{
In a rare case, if not found, let $\sk$ be a freshly sampled \ppps key subject
to $x \in \Set(\sk)$, 
and let $p = 0$.}
\item 
Find the first unconsumed 
replacement entries  
$(r_1, \DB[r_1])$ retrieved from right server,
such that $r_1$ is in $\chunk(x)$.
\footnote{In a rare case, if such an $r_1$ is not found, let it be a random
index in $\chunk(x)$, and use $0$ whenever $\DB[r_1]$ is needed later.} 
                        \item $(\sk_1',j_1)\gets\Program(\sk,\chunk(x),r_1\text{ mod } \sqrt{n})$.
                        \item 
                        Send $\sk'_1$ to the left server.     
                \end{itemize}

                \item
               {\bf Step 2: (Client Reconstructing)}
               \begin{itemize}[leftmargin=5mm,itemsep=1pt]
                   \item Receive $(\beta_{0,1},\ldots,\beta_{n^{1/4}-1,1})$ from the left server.
                   \item Save the answer as $y=p\xor  \DB[r_1]\xor\beta_{j_1,1}$.
               \end{itemize}
               
                \item
               {\bf Step 3: (Client Refreshing)}
               \begin{itemize}[leftmargin=5mm,itemsep=1pt]
                   
                   \item Sample $\sk_2$ such that $x\in\Set(\sk_2)$.
                   \item 
                   Find the first unconsumed 
                   replacement entries  
                   $(r_2, \DB[r_2])$ retrieved from left server, 
                   such that $r_2$ is in $\chunk(x)$.
                   \item $(\sk_2',j_2)\gets\Program(\sk_2,\chunk(x),r_2\text{ mod } \sqrt{n})$. 
                   \item Send $\sk_2'$ to right server.
                   \item Receive $(\beta_{0,2},\ldots,\beta_{n^{1/4}-1,2})$
                   from the right server.        
                   \item Replace the hint $(\sk,p)$ with $(\sk_2,\DB[r_2]\xor \beta_{j_2,2}\xor y)$ in the table.
               \end{itemize}
               
               \item 
               {\bf Server Responding: (Same for Left and Right Server)}
               \begin{itemize}[leftmargin=5mm,itemsep=1pt]
                   \item  Upon receiving $\sk'$, compute $(S_0,S_1,\ldots,S_{n^{1/4}-1})\gets \ListDecode(\sk')$.
                   \item Return $(\beta_0,\dots,\beta_{n^{1/4}-1})$ to the client where $\beta_i=\oplus_{i\in S_b}\DB[i]$.
               \end{itemize} %\hfill {\it // $L= n^{1/4}$}
               
            \end{enumerate}
        \end{mdframed}
    \end{minipage}
    \caption{Two-server preprocessing PIR with ${O}_\lambda(n^{1/4})$ communication, ${O}_\lambda(n^{1/2})$ computation based on PRFs. 
    }
    \label{fig:2server-new}
\end{figure}





\subsection{Construction}

\paragraph{Intuition.} 
The scheme has three major components.
\begin{itemize}[leftmargin=*]
    \item \textit{Preprocessing.} The client randomly samples $\tilde{O}(\sqrt{n})$ privately programmable pseudorandom sets, each of size $\sqrt{n}$. 
    It queries the \underline{right server} for the sets' parities, storing them along with the keys.
    Moreover, the client queries the \underline{right server} for the values of logarithmic numbers of randomly sampled indices for each $\sqrt{n}$-size chunk.
    Those entries are stored as the ``replacement'' entries.
    
    \item \textit{Online Query.} 
    Given a query $x$, the client finds a set $S$ such that $x\in S$.
    The client then finds a replacement entry $r$ that resides in the same chunk as $x$. 
    The client privately programs the set, intending to change it from $S$ to $(S/\{x\})\cup \{r\}$.
    Once the client knows the parity for this new set, it can compute $\DB[x]$ because it already knows $\DB[r]$ and the parity for $S$.
    The client uses the PPPS programming function to program the set, and sends the programmed key $\sk'$ to the \underline{left server}.
    The \textit{left server} runs the list decoding algorithm, then computes and returns all $n^{1/4}$ candidate sets' parities.
    The client knows that there is one candidate parity corresponding to the correct set $(S/\{x\})\cup \{r\}$, which is enough to compute the answer.
    
    \item \textit{Refresh.} 
    Each query consumes a set. %The client needs to replace the consumed set with a new set.
    After each query, the client just samples a new set conditioned on it containing the query $x$, and queries the \textit{right server} for its parity with the same query technique mentioned above. 
    The new set will replace the consumed set.
\end{itemize}

\ignore{
\paragraph{Intuition of the PIR scheme based on PPPS.} 
During the preprocessing phase, the client randomly samples $M_1=\tilde{O}(\sqrt{n})$ privately programmable pseudorandom sets, whose keys are denoted as $\sk_1\dots,\sk_{M_1}$, and queries the \textit{right server} for the parities for all these sets, denoted as $p_1,\dots,p_{M_1}$.
Moreover, the client queries the \textit{right server} for $M_2=\tilde{O}(1)$ random ``replacement'' entries for each $\sqrt{n}$-size chunk, and stores their indices and values as the form of $(r,\DB[r])$.
During the online phase, given a query $x$, the client finds a key $\sk$ that $\Set(\sk)$ contains $x$. The set's parity is denoted as $p$.
The client then finds a replacement entry $r$ that is in the same chunk as $x$, and it privately programs the set, intending to change it from $\Set(\sk)$ to $(\Set(\sk)/\{x\})\cup \{r\}$.
Once the client knows the parity for this new set, it can computes $\DB[x]$ because it already knows $\DB[r]$ and the parity for $\Set(\sk)$.
Specifically, the client programs the offset of $r$ in $\chunk(x)$ to the key $\sk$ using the PPPS programming function, and sends the programmed key $\sk'$ to the \textit{left server}.
The \textit{left server} runs the list decoding algorithm, then computes and returns all $n^{1/4}$ candidate sets' parities.
The client knows that there is one candidate parity corresponding to the correct set $(\Set(\sk)/\{x\})\cup \{r\}$ and it is enough to compute the answer.
To support multiple queries, the client needs to replace the consumed key with a new one.
After each query, the client just samples a new PPPS key conditioned on its set containing the query $x$, and queries the \textit{right server} for the parity for this new set with the same private query technique mentioned above. 
The client will replace the consumed key with this new key.
}

\paragraph{Detailed algorithm for bounded, random queries.}
We describe the detailed construction
for $Q = \sqrt{n} \log \kappa \cdot \alpha$ random, distinct queries 
in \Cref{fig:2server-new}.
We can easily extend such a scheme to support
unbounded, arbitrary queries using known techniques~\cite{pir-ec23}. 
For completeness, we explain how the extension works shortly after.



\paragraph{Efficiency.}
Observe that the list decoding
produces $n^{1/4}$ candidate sets
each of size $\sqrt{n}$.
Naively, expanding all sets and computing their corresponding parities takes $O_\lambda(n^{3/4})$ time.
We are still going to rely on the fact that $\ListDecode$ has an $O(\sqrt{n})$-size succinct representation to optimize the computation. 
Recall that we can first compute the common set of size $\sqrt{n}$, and then the symmetric difference between each possible decoding set and the common set will only contain $2n^{1/4}$ elements. 
Therefore, to compute the parities for all $n^{1/4}$ possible sets, we first compute the parity for the common set $S$, which takes $O_\lambda(\sqrt{n})$ time.
Then, it takes $O_\lambda(n^{1/4})$ time to enumerate the symmetric difference between the $i$-th set and the common set, which suffices to compute the parity for the $i$-th set.
So the total computation time will be $O_\lambda(\sqrt{n})$.

\paragraph{Supporting unbounded, arbitrary queries.}
For completeness, we review the techniques described in previous works 
for upgrading the scheme for $Q$ random, distinct queries to a scheme
supporting unbounded, arbitrary queries.
We can easily get rid of the distinct query
assumption in the following way: we require the client to store a local cache of size $Q$ for caching the
most recent $Q$ queries. If the client wants a repeated query, it can lookup
in the cache and make a distinct fake query. 

Further, we can assume that the queries are random without loss of generality as follows: 
we can let the client and the servers agree on 
a small-domain pseudorandom permutation (PRP)~\cite{prp1,prp2} 
(which is implied by one-way functions~\cite{swapornot})
upfront and the server can permute the database according to the PRP. 
Another option is let one of the servers
build the database as a key-value storage and use a cuckoo hash
table~\cite{cuckoo1,cuckoo2} directly based on a PRF to locate the queries, and share it with the other server.
Notice that, in both implementations, the client can still make queries
adaptively depending on the real query sequence and the responses, which is
sufficient for practical usage. Then, as long as the client makes the
queries independent of the randomness of the PRP/PRF, those queries can be
considered as uniformly random. This assumption is only needed for the
correctness.

Lastly, we can remove the bounded $Q$ query assumption as follows: 
we use a pipelining trick suggested in earlier works~\cite{pir-ec23,piano}. 
Essentially, we can spread the pre-processing for the next window of $Q$ queries
over the current window of $Q$ queries.
\ignore{
The scheme first runs a preprocessing phase for the first $Q$ queries. 
Then, during the query phase, the client interacts with the preprocessing server and gradually finishes the preprocessing for the next $Q$ queries. 
Essentially, the client and the server can split the work (including computation and communication) of the preprocessing into $Q$ pieces.
Whenever the client finishes a query, the client and the preprocessing server will finish the next piece of work about preprocessing.
Once the client finishes the $Q$ queries, it can immediately start the next window of $Q$ queries because the preprocessing for the next window is done. 
As stated above, the preprocessing phase needs ${O}_\lambda(\sqrt{n})$ communication and ${O}_\lambda(n)$ in total.  Thus, in our final scheme, we have a one-time preprocessing phase with ${O}_\lambda(n)$ communication, ${O}_\lambda(n)$ computation, while consuming only ${O}(\sqrt{n})$ client storage. After the one-time preprocessing,
we can support an unbounded number of arbitrary queries. The communication and computation
cost of each query is ${O}_\lambda({n}^{1/4})$ and the computation cost of each query is ${O}_\lambda(\sqrt{n})$.
}






\begin{theorem}
    \label{thm:eff2srv} %Let $\lambda,\kappa\in \mathbb{N}$.
    Let $\alpha(\kappa)$ be any superconstant function. 
Suppose that $\PRF_1,\PRF_2$ are secure pseudorandom functions, and that 
$n$ is bounded by $\poly(\lambda)$ and
    $\poly(\kappa)$. The two-server scheme in~\Cref{fig:2server-new} 
that supports
    $Q = \sqrt{n}\log \kappa \cdot \alpha$ 
random, distinct
queries is private, and correct with probability
    $1-\negl(\lambda)-\negl(\kappa)$ for some negligible function 
$\negl(.)$. Further, it achieves the following performance
    bounds:
    \begin{itemize}
        \item $O_\lambda(\sqrt{n}\log \kappa \alpha(\kappa))$ client storage and
        no additional server storage;
        \item \textbf{Pre-processing Phase}:
        \begin{itemize}
            \item $O_\lambda(n \log \kappa \cdot \alpha)$ 
server time and $O_\lambda(\sqrt{n}\log \kappa \cdot \alpha)$ client time;
            \item $O_\lambda(\sqrt{n} \log \kappa \cdot \alpha)$ communication;
        \end{itemize}
        \item \textbf{Query Phase}:
        \begin{itemize}
            \item 
$O_\lambda(\sqrt{n})$ expected client time 
and $O_\lambda(\sqrt{n})$ server
            time per query; 
            \item $O_\lambda(n^{1/4})$
            communication per query.
        \end{itemize}
    \end{itemize} 
Therefore, the amortized communication per query
is %\elaine{FILL}  
$O_\lambda(n^{1/4})$, and the amortized server computation
and expected client computation
per query is %$O_\lambda(\sqrt{n}\log \kappa \cdot \alpha)$, 
$O_\lambda(\sqrt{n})$.
\end{theorem}
\begin{proof}
We defer the privacy and correctness proofs to \Cref{sec:2server-priv}
and \Cref{sec:2server-correct}
respectively. Here, we focus on proving the efficiency claims.

The client stores $M_1 = \sqrt{n} \log \kappa \cdot \alpha$ number
of \ppps keys, and $M_2 = 3 \log \kappa \cdot \alpha$ number
of replacement entries per chunk.
Therefore, the space required is $O_\lambda(\sqrt{n} \log \kappa \cdot \alpha)$.

During the offline phase, the client
sends $M_1$ \ppps  
keys to the right server, and sends $M_2$ indices per chunk
to either server for constructing replacement entries.
Therefore, the offline communication is bounded
by $O_\lambda(\sqrt{n} \log \kappa \cdot \alpha)$. 
The right server needs to expand the sets for each \ppps key received
and evaluate the xor-sums.  
Both servers need to return the bits for the replacement entries.
Therefore, the total server computation is bounded
by $O_\lambda(n \log\kappa \cdot \alpha)$.

During the query phase, the client
sends one programmed \ppps key to each server,
and the size of a programmed key is at most $O_\lambda(n^{1/4})$.
Each server sends back the xor-sums of $n^{1/4}$
candidate sets.
Each candidate set has size $n^{1/2}$, however, all
$n^{1/4}$ candidate sets has a succinct representation
of size only $n^{1/2}$, 
and server can compute this succinct representation
in time $O_\lambda(n^{1/2})$.
Further, it is not hard to see that due to the structure
of the candidate sets, the server can compute
all $n^{1/4}$ xor-sums in time only $O(n^{1/2})$.
Therefore, the servers' running time is
bounded by $O_\lambda(n^{1/2})$
during each query query.
The client needs to find a matched hint, and compute
$O(1)$ xor operations during each query.
Its running time is 
dominated by the cost of finding a matched hint, which
can be done by invoking the set membership operation
for each of the $M_1$ hints.
Using \Cref{lem:2server-hintfresh}
and the pseudorandomness property of the \ppps, the expected
number of hints checked until a key $\sk$ such that $\Set(\sk)$ contains the current query is found is $O(\sqrt{n})$.
The expected number of tries till success is $\sqrt{n}$.
Therefore, the client's expected running time per query
is upper bounded by $O_\lambda(\sqrt{n})$.
\end{proof}



\subsection{Privacy Proof}
\label{sec:2server-priv}

Suppose that the underlying \ppps scheme satisfies
private programmability.
Below, we prove the privacy of our two-server PIR scheme.

In the pre-processing 
phase, the server sends the sets $\Set(\sk_i)$ only to the right server, thereby no information about these sets is leaked to the left server. Similarly, no information about the indices $r_1$ is leaked to the right server and no information about the indices $r_2$ is leaked to the left server.
We will first the lemma about the distribution of client's hint table, when the adversary controls either of the left or right server.
\begin{lemma}
 Recall that in each time step $t$, the adversary \algA 
adaptively chooses a query $x_t \in \{0, 1, \ldots, n-1\}$ for the client.
%The following statements are true:
%\begin{enumerate}[leftmargin=5mm]
%\item 
At the end of each time step $t$, 
the client's hint table  
is distributed 
as a table of size $M_1$
where each entry is a freshly sampled \ppps key, even when
conditioned on $\algA$'s view so far.  
%\end{enumerate}
\label{lem:2server-hintfresh}
\end{lemma}
\begin{proof}
Suppose the above statement holds at the end of time step $t-1$.
We prove that it still holds 
at the end of time step $t$.
Since the hint table is distributed as a fresh randomly sampled
table even when conditioned on $\algA$'s view
at the end of $t-1$, we may henceforth assume
an arbitrary fixed query $x_t$.
The distribution of the hint table 
before the $t$-th query can be equivalently rewritten as:
\begin{itemize}[leftmargin=5mm]
\item First, 
sample the decision whether any 
of the $M_1$ entries contains the current query $x_t$,
and if so, which is the first entry (denoted $i^*$) 
that contains $x_t$. 
If not found, we assume $i^* = M_1 + 1$. 
\item 
For each $i < i^*$, sample a random \ppps
key subject to not containing $x_t$.
\item 
For each $i = i^*$, sample a random \ppps
key subject to containing $x_t$.
\item 
For each $i > i^*$, sample a random \ppps key.
\end{itemize}

Using the above interpretation, it is easy to see 
that the distribution of the hint table 
after the $t$-th query is unaltered 
no matter which of the two servers \algA controls.
\end{proof}


\subsubsection{Left Server Privacy}
We first construct 
the following simulator for proving left server privacy.

\paragraph{Simulator construction.}
\begin{itemize}[leftmargin=5mm]
\item  
During the pre-processing phase, 
for each chunk $\ell$, sample $M_2$ 
random indices belonging to $\ell$, send them to $\algA$.
\item 
During each query, 
call the simulator of the \ppps scheme which outputs $\sk'$, 
send $\sk'$ to $\algA$. 
\end{itemize}

\paragraph{Indistinguishability of ${\sf Real}$ and ${\sf Ideal}$.}
We now prove the indistinguishability of the ${\sf Real}$ and
${\sf Ideal}$ for both the servers
assuming the 
private programmability of the underlying \ppps scheme.

First, due to \Cref{lem:2server-hintfresh}, 
we can equivalently rewrite the ${\sf Real}$ experiment 
for the right server as follows:
at the end of each time step, resample the entire hint table
freshly at random before continuing to answer more queries.
As a result, 
the view of $\algA$ who controls the right server is 
distributed as: 
\begin{itemize}[leftmargin=5mm]
\item {\it Pre-processing phase}.
%Sample $M_1$ random \ppps keys, and send them to $\algA$. 
For each chunk $\ell$, send $M_2$ random indices
in chunk $\ell$ to \algA.
\item {\it Each time step $t$}.
\begin{itemize}[leftmargin=5mm]
\item 
sample a \ppps key $\sk$
 at random subject to containing the query $x_t$;
sample $\delta$ at random
from $\{0, \ldots, \sqrt{n}-1\}$.
\item 
call $\sk', \_ \leftarrow {\sf Program}(\sk, \chunk(x_t), \delta)$; 
\item 
send $\sk'$ to \algA.
\end{itemize}
\end{itemize}
One way to see this is to think of the distribution of the table
as the equivalent distribution
in the proof of \Cref{lem:hintfresh}. 
Further, observe that each $r_2\text{ mod }\sqrt{n}$ in the scheme 
is distributed randomly from the perspective
of the left server,  
since they 
were only sent to the 
right server during the pre-processing phase.
%because the right server has no information about $r_1$ during preprocessing.
%For the right server, the distribution of $\sk'$ is the actual distribution and $r_1\text{ mod }\sqrt{n}$ is distributed randomly because the right server has no information about $r_2$ during preprocessing.

Therefore, the rest of the proof follows
due to a straightforward hybrid argument
where we replace
the programmed keys (denoted $\sk'$ earlier) 
sent to the right server in all time steps one by one
with a simulated key, 
relying on the private programmability of the underlying \ppps. 






\subsubsection{Right Server Privacy}

We first construct 
the following simulator for proving right server privacy.

\paragraph{Simulator construction.}
\begin{itemize}[leftmargin=5mm]
\item  
During the pre-processing phase, 
send $M_1$ randomly sampled \ppps keys to $\algA$.
Further, for each chunk $\ell$, sample $M_2$ 
random indices in $\ell$, send them to $\algA$.
\item 
During each query, 
call the simulator of the \ppps scheme which outputs $\sk'$, 
%sample a random \ppps key $\sk$ subject to containing $x_t$, 
%sample a random $\delta \in \{0, \ldots, \sqrt{n}-1\}$, 
%call $\sk',\_ \leftarrow {\sf Program}(\sk, 
%\chunk(x_t), \delta)$, 
send $\sk'$ to $\algA$. 
\end{itemize}


\paragraph{Indistinguishability of ${\sf Real}$ and ${\sf Ideal}$.}
We now prove the indistinguishability of the ${\sf Real}$ and
${\sf Ideal}$ for both the servers
assuming the 
private programmability of the underlying \ppps scheme.
\ignore{
First, due to \Cref{lem:2server-hintfresh}, 
we can equivalently rewrite the ${\sf Real}$ experiment 
for the right server as follows:
at the end of each time step, resample the entire hint table
freshly at random before continuing to answer more queries.
As a result, 
}
The view of $\algA$ who controls the right server is 
distributed as: 
\begin{itemize}[leftmargin=5mm]
\item {\it Pre-processing phase}.
Sample $M_1$ random \ppps keys, and send them to $\algA$. 
Further, for each chunk $\ell$, send $M_2$ random indices
in chunk $\ell$ to \algA.
\item {\it Each time step $t$}.
\begin{itemize}[leftmargin=5mm]
\item 
sample a \ppps key $\sk$
 at random subject to containing the query $x_t$;
sample $\delta$ at random
from $\{0, \ldots, \sqrt{n}-1\}$.
\item 
call $\sk', \_ \leftarrow {\sf Program}(\sk, \chunk(x_t), \delta)$; 
\item 
send $\sk'$ to \algA.
\end{itemize}
\end{itemize}
%One way to see this is to think of the distribution of the table
%as the equivalent distribution
%in the proof of \Cref{lem:hintfresh}. 
To see the above, observe that each 
$r_1\text{ mod }\sqrt{n}$ in the scheme 
is distributed randomly from the perspective
of the right server,  
since they 
were only sent to the 
left server during the pre-processing phase.
%because the right server has no information about $r_1$ during preprocessing.
%For the right server, the distribution of $\sk'$ is the actual distribution and $r_1\text{ mod }\sqrt{n}$ is distributed randomly because the right server has no information about $r_2$ during preprocessing.

Therefore, the rest of the proof follows
due to a straightforward hybrid argument
where we replace
the programmed keys (denoted $\sk'$ earlier) 
sent to the right server in all time steps one by one
with a simulated key, 
relying on the private programmability of the underlying \ppps. 




\ignore{
\paragraph{Simulator construction and the ${\sf Ideal}$ experiment.}
Consider the following simulator construction for both the left server which does not make
use of the queries: 
in every time step $t$, call the simulator $\Sim$ of the \ppps
scheme, and let the output be $\sk'$. 
Send $\sk'$ to the server. 

\paragraph{Indistinguishability of ${\sf Real}$ and 
${\sf Ideal}$.}
We now prove the indistinguishability of the ${\sf Real}$ and
${\sf Ideal}$ for both the servers
assuming the 
private programmability of the underlying \ppps scheme.

First, due to \Cref{lem:hintfresh}, 
we can equivalently rewrite the ${\sf Real}$ experiment for the left and te right server as follows:
at the end of each time step, resample the entire hint table
freshly at random before continuing to answer more queries.
As a result, the %distribution of the $\sk'$ received by 
messages sent to 
\algA in each time step $t$
is distributed as 
\begin{itemize}[leftmargin=5mm]
\item[]
\begin{itemize}[leftmargin=5mm]
\item 
sample a \ppps key $\sk$
 at random subject to containing the query $x_t$;
sample $\delta$ at random
from $\{0, \ldots, \sqrt{n}-1\}$.
\item 
call $\sk', \_ \leftarrow {\sf Program}(\sk, \chunk(x_t), \delta)$; 
\item 
send $\sk'$ to \algA.
\end{itemize}
\end{itemize}
One way to see this for the left server is to think of the distribution of the table
as the equivalent distribution
in the proof of \Cref{lem:hintfresh}. Moreover, $r_1\text{ mod }\sqrt{n}$ is distributed randomly because the right server has no information about $r_1$ during preprocessing.
For the right server, the distribution of $\sk'$ is the actual distribution and $r_1\text{ mod }\sqrt{n}$ is distributed randomly because the right server has no information about $r_2$ during preprocessing.

Therefore, the rest of the proof follows
due to a straightforward hybrid argument
where we replace
the programmed keys (denoted $\sk'$ earlier) 
sent to the server in all time steps one by one
with a simulated key, 
relying on the private programmability of the underlying \ppps. 
}


\subsection{Correctness Proof}
\label{sec:2server-correct}

We show that with $Q = \sqrt{n} \log \kappa \cdot \alpha$ 
random, distinct queries,
the probability of ever having correctness error is negligibly small. 
An error can happen 
if one of the following bad events take place:
\begin{itemize}[leftmargin=5mm]
\item 
{\it No matched hint}. 
During some query for $x$, 
no hint is found that contains the query $x$.
\item 
{\it Depleting replacement entries}. 
During some query for $x$, there is no more replacement
entry of the form $(r_1, \DB[r_1])$ 
or $(r_2, \DB[r_2])$ 
corresponding to $\chunk(x)$. 
\end{itemize}

Below, we show 
that the probability of each bad event during a window
of $Q$ random, distinct queries 
is negligibly small. 

\paragraph{Probability of no matched hint.}
Due to \Cref{lem:2server-hintfresh}, for any fixed time step $t$, 
we can assume the client's hint table
contains freshly sampled \ppps keys and is independent of 
the current query $x_t$.
Due to the pseudorandomness property of the \ppps,
the sets generated by the keys 
in the hint table are computationally indisintinguishable
from 
$M_1$ sets independently sampled from the distribution $\mcal{D}_n$.
Below we 
calculate the probability that a fixed  
element $x_t$ is not in any of the $M_1$ sets 
sampled independently from 
$\mcal{D}_n$ --- the probability that $x_t$ is not contained in any entry in the
client's hint table can only be negligibly different.

The probability that one set 
sampled from $\mcal{D}_n$ contains $x_t$ is $1/\sqrt{n}$. 
Therefore, the probability that 
none of the $M_1$ sets contains $x_t$ is
$(1-1/\sqrt{n})^{M_1}$, and given the choice
of $M_1 = \sqrt{n} \log \kappa \alpha(\kappa)$ where $\alpha(\kappa)$
is a super-constant function,
this probability is negligibly small in $\kappa$.

Finally, taking a union bound over all polynomially many time steps,
the probability of ever not having a matched hint
is negligibly small in $\kappa$. 


\paragraph{Probability of depleting replacement entries.} 
One can only deplete the replacement 
entries of some chunk $\ell$ if the chunk $\ell$
is encountered more than $M_2$ times.
With $Q$ random distinct queries, 
each query will hit a random chunk.
The expected number of hits per chunk
is therefore $Q/\sqrt{n} = \log \kappa \cdot \alpha$.
By the Chernoff bound, 
the probability that the number of visits to some fixed chunk $\ell$
exceeds $M_2 = 3 \log \kappa \cdot \alpha$ is negligibly small
in $\kappa$ as long as $\alpha(\kappa)$ is a super-constant function.

Finally, taking a union bound over 
all chunks and 
all polynomially many time steps,
the probability of ever 
depleting replacement entries of any chunk
is negligibly small in $\kappa$. 



